\documentclass[12pt,a4paper]{article}
\usepackage{amsmath, amssymb}
\usepackage[margin=2.5cm]{geometry}
\usepackage{enumitem}
\usepackage{xcolor}

\newcommand{\easy}{\textcolor{green!60!black}{$\bigstar$}}
\newcommand{\medium}{\textcolor{orange}{$\bigstar\bigstar$}}
\newcommand{\hard}{\textcolor{red}{$\bigstar\bigstar\bigstar$}}

\title{\textbf{Solutions: Additional Exercises 1}\\Positive Series}
\author{Series and Integral Transforms}
\date{}

\begin{document}
\maketitle

% -------------------------------------------------------
\section*{Exercise 1 — Finding Exact Sums}

\subsection*{(a) \easy\ $\displaystyle\sum_{n=1}^{\infty}\left(\frac{3}{4}\right)^n$}

This is a \textbf{geometric series} with first term $a = \tfrac{3}{4}$ and ratio $r = \tfrac{3}{4}$.
Since $|r| < 1$, the series converges and
\[
  \sum_{n=1}^{\infty}\left(\frac{3}{4}\right)^n
  = \frac{a}{1-r}
  = \frac{\tfrac{3}{4}}{1 - \tfrac{3}{4}}
  = \frac{\tfrac{3}{4}}{\tfrac{1}{4}} = \boxed{3}.
\]

\subsection*{(b) \easy\ $\displaystyle\sum_{n=0}^{\infty}\frac{2^n+3^n}{5^n}$}

Split the series using linearity:
\[
  \sum_{n=0}^{\infty}\frac{2^n+3^n}{5^n}
  = \sum_{n=0}^{\infty}\left(\frac{2}{5}\right)^n + \sum_{n=0}^{\infty}\left(\frac{3}{5}\right)^n.
\]
Both are geometric series (starting at $n=0$) with $|r|<1$:
\[
  = \frac{1}{1-\tfrac{2}{5}} + \frac{1}{1-\tfrac{3}{5}}
  = \frac{5}{3} + \frac{5}{2}
  = \frac{10}{6}+\frac{15}{6} = \boxed{\frac{25}{6}}.
\]

\subsection*{(c) \easy\ $\displaystyle\sum_{n=1}^{\infty}\frac{1}{n(n+1)}$}

Use \textbf{partial fractions}: $\dfrac{1}{n(n+1)} = \dfrac{1}{n} - \dfrac{1}{n+1}$.

The partial sum \textbf{telescopes}:
\[
  S_N = \sum_{n=1}^{N}\left(\frac{1}{n}-\frac{1}{n+1}\right)
  = 1 - \frac{1}{N+1} \;\xrightarrow{N\to\infty}\; \boxed{1}.
\]

\subsection*{(d) \easy\ $\displaystyle\sum_{n=1}^{\infty}\frac{1}{n(n+2)}$}

Partial fractions: $\dfrac{1}{n(n+2)} = \dfrac{1}{2}\!\left(\dfrac{1}{n}-\dfrac{1}{n+2}\right)$.

Telescoping with step 2:
\[
  S_N = \frac{1}{2}\!\left(1 + \frac{1}{2} - \frac{1}{N+1} - \frac{1}{N+2}\right)
  \;\xrightarrow{N\to\infty}\;
  \frac{1}{2}\cdot\frac{3}{2} = \boxed{\frac{3}{4}}.
\]

\subsection*{(e) \medium\ $\displaystyle\sum_{n=1}^{\infty}\!\left[\sin\!\left(\frac{1}{n}\right)-\sin\!\left(\frac{1}{n+1}\right)\right]$}

Define $b_n = \sin(1/n)$. The general term is $b_n - b_{n+1}$, which \textbf{telescopes}:
\[
  S_N = b_1 - b_{N+1} = \sin(1) - \sin\!\left(\frac{1}{N+1}\right)
  \;\xrightarrow{N\to\infty}\;
  \sin(1) - 0 = \boxed{\sin 1}.
\]

\subsection*{(f) \medium\ $\displaystyle\sum_{n=2}^{\infty}\ln\!\left(1-\frac{1}{n^2}\right)$}

Factor the argument: $1 - \dfrac{1}{n^2} = \dfrac{(n-1)(n+1)}{n^2}$.  Using log rules:
\[
  \ln\!\left(\frac{(n-1)(n+1)}{n^2}\right)
  = \ln(n-1) - \ln n + \ln(n+1) - \ln n.
\]
Write $a_n = \ln(n-1) - \ln n$ and $c_n = \ln(n+1) - \ln n$. Then:
\[
  S_N = \sum_{n=2}^{N} a_n + \sum_{n=2}^{N} c_n.
\]
The first sum telescopes to $\ln 1 - \ln N = -\ln N$.  The second telescopes to $\ln(N+1) - \ln 2$.  Hence
\[
  S_N = -\ln N + \ln(N+1) - \ln 2
  = \ln\!\left(\frac{N+1}{N}\right) - \ln 2
  \;\xrightarrow{N\to\infty}\; 0 - \ln 2 = \boxed{-\ln 2}.
\]

\subsection*{(g) \easy\ $\displaystyle\sum_{n=1}^{\infty}\frac{4^{n+1}}{5^n}$}

Write $\dfrac{4^{n+1}}{5^n} = 4\cdot\left(\dfrac{4}{5}\right)^n$.  Geometric series:
\[
  4\sum_{n=1}^{\infty}\left(\frac{4}{5}\right)^n
  = 4\cdot\frac{\tfrac{4}{5}}{1-\tfrac{4}{5}}
  = 4\cdot\frac{\tfrac{4}{5}}{\tfrac{1}{5}} = 4\cdot 4 = \boxed{16}.
\]

\subsection*{(h) \medium\ $\displaystyle\sum_{n=0}^{\infty}\frac{1}{(3n+1)(3n+4)}$}

Partial fractions: $\dfrac{1}{(3n+1)(3n+4)} = \dfrac{1}{3}\!\left(\dfrac{1}{3n+1}-\dfrac{1}{3n+4}\right)$.

Telescoping (the terms pair up with step 3):
\[
  S_N = \frac{1}{3}\!\left(\frac{1}{1} - \frac{1}{3N+4}\right)
  \;\xrightarrow{N\to\infty}\;
  \frac{1}{3}\cdot 1 = \boxed{\frac{1}{3}\cdot\frac{1}{1}} .
\]
Wait — at $n=0$ the first term is $\frac{1}{3}(\frac{1}{1}-\frac{1}{4})$, at $n=1$ it is $\frac{1}{3}(\frac{1}{4}-\frac{1}{7})$, etc.  The telescope collapses to
\[
  S_N = \frac{1}{3}\left(1 - \frac{1}{3N+4}\right) \to \frac{1}{3}.
\]
\textbf{But} note the first factor at $n=0$ is $3(0)+1=1$ and $3(0)+4=4$, so indeed $S = \boxed{\dfrac{1}{3}}$.

\smallskip
\textit{Correction check:} At $n=0$: $1/(1\cdot4)=1/4$.  At $n=1$: $1/(4\cdot7)$.  The telescope gives $S=\frac{1}{3}(1/1)=1/3$.  $\checkmark$

% -------------------------------------------------------
\section*{Exercise 2 — Convergence Tests}

\subsection*{(a) \easy\ $\displaystyle\sum\frac{1}{n^2+1}$}

Since $\dfrac{1}{n^2+1} \leq \dfrac{1}{n^2}$ and $\sum\dfrac{1}{n^2}$ converges (\textbf{$p$-series}, $p=2>1$), by the \textbf{Comparison Test} the series \textbf{converges}.

\subsection*{(b) \easy\ $\displaystyle\sum\frac{n}{2^n}$}

\textbf{Ratio Test}: $\dfrac{a_{n+1}}{a_n} = \dfrac{n+1}{2^{n+1}}\cdot\dfrac{2^n}{n} = \dfrac{n+1}{2n}\to\dfrac{1}{2}<1$. \textbf{Converges}.

\subsection*{(c) \easy\ $\displaystyle\sum\frac{n^2}{n!}$}

\textbf{Ratio Test}: $\dfrac{(n+1)^2}{(n+1)!}\cdot\dfrac{n!}{n^2} = \dfrac{(n+1)^2}{(n+1)\,n^2} = \dfrac{n+1}{n^2}\to 0 < 1$. \textbf{Converges}.

\subsection*{(d) \medium\ $\displaystyle\sum\frac{1}{n(\ln n)^2}$, $n\geq 2$}

\textbf{Integral Test}: $f(x)=\dfrac{1}{x(\ln x)^2}$ is positive and decreasing.
\[
  \int_2^{\infty}\frac{dx}{x(\ln x)^2}
  = \left[-\frac{1}{\ln x}\right]_2^{\infty} = \frac{1}{\ln 2} < \infty.
\]
\textbf{Converges}.

\subsection*{(e) \easy\ $\displaystyle\sum\frac{2^{1/n}}{3^n}$}

Since $2^{1/n}\to 1$, we have $\dfrac{2^{1/n}}{3^n}\sim\dfrac{1}{3^n}$ as $n\to\infty$.  By \textbf{Limit Comparison} with $\sum(1/3)^n$ (geometric, converges): $L = \lim \dfrac{2^{1/n}/3^n}{1/3^n} = 1$.  \textbf{Converges}.

\subsection*{(f) \easy\ $\displaystyle\sum\frac{\arctan(1/n)}{n}$}

Since $\arctan(1/n)\sim 1/n$ as $n\to\infty$, the terms behave like $1/n^2$.  \textbf{Limit Comparison} with $\sum 1/n^2$: $L=\lim n^2\cdot\arctan(1/n)/n = \lim n\arctan(1/n)=1$.  \textbf{Converges}.

\subsection*{(g) \hard\ $\displaystyle\sum\frac{1\cdot3\cdot5\cdots(2n-1)}{(2n)!}$}

\textbf{Ratio Test}:
\[
  \frac{a_{n+1}}{a_n}
  = \frac{1\cdot3\cdots(2n-1)(2n+1)}{(2n+2)!}\cdot\frac{(2n)!}{1\cdot3\cdots(2n-1)}
  = \frac{2n+1}{(2n+2)(2n+1)} = \frac{1}{2n+2}\to 0 < 1.
\]
\textbf{Converges}.

\subsection*{(h) \medium\ $\displaystyle\sum\left(\frac{n}{3n-1}\right)^{2n-1}$}

\textbf{Root Test}: $a_n^{1/n} = \left(\dfrac{n}{3n-1}\right)^{(2n-1)/n} \to \left(\dfrac{1}{3}\right)^2 = \dfrac{1}{9} < 1$. \textbf{Converges}.

\subsection*{(i) \easy\ $\displaystyle\sum\left(\frac{n+1}{2n-1}\right)^n$}

\textbf{Root Test}: $a_n^{1/n} = \dfrac{n+1}{2n-1}\to\dfrac{1}{2}<1$. \textbf{Converges}.

\subsection*{(j) \medium\ $\displaystyle\sum\frac{n!}{1\cdot3\cdot5\cdots(2n-1)}$}

\textbf{Ratio Test}:
\[
  \frac{a_{n+1}}{a_n} = \frac{(n+1)!}{1\cdot3\cdots(2n+1)}\cdot\frac{1\cdot3\cdots(2n-1)}{n!}
  = \frac{n+1}{2n+1}\to\frac{1}{2}<1.
\]
\textbf{Converges}.

\subsection*{(k) \hard\ $\displaystyle\sum\frac{(2n)!}{4^n(n!)^2}$}

\textbf{Ratio Test}: $\dfrac{a_{n+1}}{a_n} = \dfrac{(2n+2)!}{4^{n+1}((n+1)!)^2}\cdot\dfrac{4^n(n!)^2}{(2n)!} = \dfrac{(2n+2)(2n+1)}{4(n+1)^2} = \dfrac{2n+1}{2(n+1)}\to 1$. \textit{Inconclusive.}

By \textbf{Stirling's approximation}, $(2n)!\sim\sqrt{4\pi n}(2n/e)^{2n}$ and $n!\sim\sqrt{2\pi n}(n/e)^n$, giving:
\[
  \frac{(2n)!}{4^n(n!)^2}\sim\frac{1}{\sqrt{\pi n}}.
\]
Since $\sum 1/\sqrt{\pi n}$ diverges ($p$-series, $p=1/2\leq 1$), the original series \textbf{diverges} by Limit Comparison.

% -------------------------------------------------------
\section*{Exercise 3 — Theoretical Questions}

\subsection*{(a) \easy\ Necessary condition}

The \textbf{Necessary Condition for Convergence} states: if $\sum a_n$ converges, then $a_n\to 0$.

\begin{itemize}
  \item $\sum \dfrac{n^2}{n^2+1}$: here $a_n = \dfrac{n^2}{n^2+1}\to 1\neq 0$. \textbf{Diverges}.
  \item $\sum (-1)^n$: here $a_n = (-1)^n$ does not converge to $0$. \textbf{Diverges}.
\end{itemize}

\subsection*{(b) \easy\ Absolute convergence and scalar multiples}

Given $\sum a_n$ absolutely convergent and $c > 0$. Then $|ca_n| = c|a_n|$, so
$\sum |ca_n| = c\sum |a_n| < \infty$. Hence $\sum ca_n$ is also absolutely convergent.

\subsection*{(c) \medium\ Divergence of the sum when one part diverges}

Suppose $\sum a_n$ converges and $\sum b_n$ diverges. We claim $\sum (a_n + b_n)$ diverges.

\textit{Proof by contradiction}: assume $\sum(a_n+b_n)$ converges. Then by linearity of convergent series,
$\sum b_n = \sum(a_n+b_n) - \sum a_n$ would also converge — a contradiction.
Therefore $\sum(a_n+b_n)$ \textbf{diverges}. $\square$

\subsection*{(d) \medium\ Bouncing ball}

A ball is dropped from height $h = 4$ m, with bounce ratio $r = 2/3$.

Total distance $= h + 2h r + 2h r^2 + \cdots = h + \dfrac{2hr}{1-r}$ (geometric, $|r|<1$):
\[
  = 4 + \frac{2\cdot4\cdot\tfrac{2}{3}}{1-\tfrac{2}{3}}
  = 4 + \frac{\tfrac{16}{3}}{\tfrac{1}{3}}
  = 4 + 16 = \boxed{20 \text{ m}}.
\]

\subsection*{(e) \medium\ $\displaystyle\sum_{n=0}^{\infty}\tan^n(x)$}

This is a geometric series with ratio $r = \tan x$.  It converges if and only if $|\tan x| < 1$, i.e.,
$x \in \left(-\dfrac{\pi}{4}+k\pi,\,\dfrac{\pi}{4}+k\pi\right)$, $k\in\mathbb{Z}$.
When it converges, the sum is:
\[
  S = \frac{1}{1-\tan x}.
\]

% -------------------------------------------------------
\section*{Exercise 4 — Integral Test Verification}

\subsection*{(a) \easy\ $\displaystyle\sum_{n=2}^{\infty}\frac{1}{n\ln n}$}

$f(x)=\dfrac{1}{x\ln x}$ is positive and decreasing on $[2,\infty)$.
\[
  \int_2^{\infty}\frac{dx}{x\ln x} = \Big[\ln(\ln x)\Big]_2^{\infty} = +\infty.
\]
By the \textbf{Integral Test}, the series \textbf{diverges}.

\subsection*{(b) \easy\ $\displaystyle\sum_{n=1}^{\infty}\frac{1}{n^2+1}$}

$f(x)=\dfrac{1}{x^2+1}$ is positive and decreasing.
\[
  \int_1^{\infty}\frac{dx}{x^2+1} = \Big[\arctan x\Big]_1^{\infty} = \frac{\pi}{2} - \frac{\pi}{4} = \frac{\pi}{4} < \infty.
\]
By the \textbf{Integral Test}, the series \textbf{converges}.

\subsection*{(c) \medium\ $\displaystyle\sum_{n=1}^{\infty}n\,e^{-n^2}$}

$f(x) = xe^{-x^2}$ is positive and decreasing on $[1,\infty)$.
\[
  \int_1^{\infty} x\,e^{-x^2}\,dx
  = \left[-\frac{1}{2}e^{-x^2}\right]_1^{\infty}
  = 0 + \frac{1}{2e} = \frac{1}{2e} < \infty.
\]
By the \textbf{Integral Test}, the series \textbf{converges}.

\bigskip
\hrule
\bigskip
\textit{Key takeaways}: Geometric series and telescoping are exact-sum workhorses.  The $p$-series criterion ($p>1$ converges), Ratio/Root tests handle factorial and exponential terms, and the Integral Test bridges series to calculus.

\end{document}
